\section{서론}

TSI는 internal transition model이 모든 state coordinate를 관계없는 수치로
취급하기보다 선언된 structural dependency를 표현해야 한다고 제안한다. Paper~3의
historical standalone 연구는 다른 simulator에서 sealed structural OOD, open-loop
rollout과 same-task criterion association을 평가한다. 이 연구들은 여기서 다루는
exact factorization을 확립하지 않고 historical Paper~4 audit의 quantity를 추정하지
않으며 서로의 replication도 아니다. 그러나 Paper~3의 새 prospective extension은
본 논문의 attribution cohort다. Root seed와 multiplicity family를 공유하며 겹치는
endpoint는 독립 evidence가 아니다.

이 benchmark의 원래 질문은 discovered structural factorization이 diagonal 또는
unstructured control보다 held-out intervention을 잘 예측하는지였다. 이 표현은 너무
강하다. Finite generator가 TSI hypothesis class 안에 있고 fitted TSI predictor는
동일한 transition equation을 사용한다. Graph와 여섯 parameter를 유일하게 회수하면
exact prediction은 분석적으로 따라온다. 따라서 primary audit는 factorization의
empirical causal effect가 아니라 exact representability와 finite-family
identifiability를 다룬다.

Structural prediction에는 raw observation에서 entity discovery, 시간에 걸친
identity, relation, structural variable과 dependency graph를 거쳐 prediction에
이르는 chain이 필요하다. 본 연구는 이 chain의 후반 단계를 다룬다. Entity set과
cross-time identity는 generator가 제공하며, 이들이 주어진 뒤의 graph recovery와
prediction을 검정한다. Raw observation에서 entity를 발견하고 그 identity를
유지하는 문제는 이 audit의 estimand 밖에 있다.

Frozen comparison은 투명한 기록으로 유지하고, 명확히 표시한 post-review
diagnostic을 추가한다. 이 확장은 cell-level average가 tie와 structural zero를 어떻게
가리는지, graph choice의 효과가 factorized head와 untied dense head에서 같은지,
그리고 실제 two-coordinate action 또는 fitted TSI family 밖 transition term에서
exactness가 어디서 실패하는지를 묻는다. 이 diagnostic은 universal model-superiority
claim을 복권하지 않고 attribution boundary를 드러낸다. 이어지는 prospective
연구는 fresh stochastic family, complete graph-by-head design, matched
seven-parameter control, learned routing, genuine two-mechanism composition,
multiplicity-adjusted world-level inference와 second-language implementation으로
확인된 design defect를 해결한다.
